In this project, we compare standard Stochastic Gradient Descent (SGD) against Robust SGD (R-SGD) under gradient uncertainty. The comparison is conducted across two distinct tasks: (i) a single-layer logistic regression on linearly separable Gaussian data, and (ii) a small two-layer neural network (MLP) on a noisy semicircles dataset. In both cases, the models are trained to minimize the mean binary cross-entropy loss without additional regularization.

\subsection{Hyperparameter Tuning and Freezing}
To ensure a fair comparison, we used a fixed mini-batch size and 
tuned the learning rate for both optimizers, as well as the robustness 
radius \(\rho\) for R-SGD (including the baseline \(\rho=0\)). 
The hyperparameter grid search was performed in a clean environment 
(without artificial gradient noise, \(\sigma=0\)), and the configuration 
that minimized the mean validation loss was selected. Once the optimal 
learning rate and \(\rho\) were found, these settings were 
frozen for the subsequent noise sweep. 

\subsection{Isotropic Gradient Noise Model}
To evaluate the robustness of the optimizers, we introduced an isotropic 
Gaussian perturbation to the per-example gradients. For a mini-batch of 
size \(b\), the observed gradient for the \(i\)-th example at optimization 
step \(k\) is modeled as:
\begin{equation}
    v_{k,i} = u_i(\theta_k) + \sigma z_{k,i}, \quad z_{k,i} \sim \mathcal{N}(0, I_p)
\end{equation}
where \(u_i(\theta_k) = \nabla_\theta \ell(t_\theta(a_i), y_i)\) is the true per-example gradient, \(\sigma \ge 0\) is the noise magnitude, and \(I_p\) is the \(p \times p\) identity matrix. The perturbations \(z_{k,i}\) are sampled independently across observations and iterations. 

Because the raw gradient magnitudes differ between the two tasks, the noise levels were chosen relative to the empirical scale of the gradients. We defined a base noise scale \(\sigma_{\text{scale}}\) (measured during a warmup phase) and evaluated the optimizers at multipliers \(\sigma \in \{0, 0.1, 0.3, 1.0, 2.5\} \cdot \sigma_{\text{scale}}\). The case \(\sigma=0\) serves as the natural baseline, representing only the inherent stochasticity of mini-batch sampling. 

\subsection{Evaluation Protocol}
For each noise level and each of the 5 evaluation seeds, we trained the models from matched initializations. To ensure a fair comparison, both SGD and R-SGD were exposed to the exact same absolute noise levels and mini-batch sequences within a given seed. The performance was evaluated on uncorrupted validation and test data, reporting the final training loss, test loss, test accuracy, runtime, and zero-step frequency.